% BHU PYQ -- Transcribed & Corrected
% Paper: AIHMN-41 / AIHMN41: Political Ideas and Institutions of Ancient India
% Exam: B.A. (Hons.) Semester IV Examination, 2025-26 (NEP)
% Category: Minor Core
% Department: Department of Ancient Indian History, Culture & Archaeology (A.I.H.C. & Archaeology)

\documentclass[11pt,a4paper]{article}
\usepackage[a4paper,top=2.5cm,bottom=2.5cm,left=2.8cm,right=2.8cm,headheight=14pt]{geometry}
\usepackage{amsmath,amssymb}
\usepackage[T1]{fontenc}
\usepackage[utf8]{inputenc}
\usepackage{lmodern,microtype,fancyhdr,enumitem,hyperref,lastpage,parskip}

\hypersetup{
    colorlinks=true,
    urlcolor=black,
    linkcolor=black,
    pdftitle={AIHMN41: Political Ideas and Institutions of Ancient India -- BHU Sem IV 2025-26 (NEP)}
}

\pagestyle{fancy}
\fancyhf{}
\renewcommand{\headrulewidth}{0.4pt}
\renewcommand{\footrulewidth}{0.4pt}
\fancyhead[L]{\small \textit{Banaras Hindu University}}
\fancyhead[R]{\small \textit{AIHMN-41: Political Ideas \& Institutions (Sem IV 2025-26)}}
\fancyfoot[C]{\small Page \thepage\ of \pageref{LastPage}}

\newlist{parts}{enumerate}{2}
\setlist[parts,1]{label=(\alph*),leftmargin=*,itemsep=4pt,topsep=2pt}

\newcommand{\pts}[1]{\hfill \textit{[#1]}}

\begin{document}

\begin{center}
  {\large\bfseries BANARAS HINDU UNIVERSITY}\\[2pt]
  {\normalsize B.A. (Hons.) Semester IV Examination, 2025-26 (NEP)}\\[6pt]
  \rule{0.8\linewidth}{0.5pt}\\[6pt]
  {\large\bfseries A.I.H.C. \& ARCHAEOLOGY}\\[4pt]
  {\large\bfseries Paper Code: AIHMN-41 : Political Ideas and Institutions of Ancient India}\\[4pt]
  {\normalsize Time: 2:30 Hours \hfill Full Marks: 70}
\end{center}

\rule{\linewidth}{0.4pt}

\smallskip

\noindent \textit{Instructions: Attempt all questions. Questions of Sections A, B and C carry 10, 05 and 02 marks each.}\\[2pt]
\noindent \textit{नोट: सभी प्रश्नों के उत्तर दीजिए। खण्ड - अ, ब एवं स के प्रश्न क्रमशः 10, 05 और 02 अंकों के हैं।}

\bigskip

\begin{center}
  {\large\bfseries SECTION A / खण्ड-अ}\\[2pt]
  {\normalsize (Long Answer Type Questions / दीर्घ उत्तरीय प्रश्न)}
\end{center}

\noindent \textit{Note: Attempt all questions. Answer each question in approximately 250 words. Each question carries 10 marks.}\\[2pt]
\noindent \textit{नोट: सभी प्रश्नों के उत्तर दीजिए। प्रत्येक प्रश्न का उत्तर लगभग 250 शब्दों में दीजिए। प्रत्येक प्रश्न 10 अंकों का है।}

\bigskip

\noindent \textbf{Question 1.} Critically examine the different theories of the origin of the State in Ancient India.\\[1pt]
प्राचीन भारत में राज्य की उत्पत्ति के विभिन्न सिद्धान्तों का समालोचनात्मक विश्लेषण कीजिए। \pts{10}

\begin{center}
  \textbf{OR / अथवा}
\end{center}

\noindent Discuss the Saptanga theory of the State and explain its components in detail.\\[1pt]
राज्य के सप्तांग सिद्धान्त की व्याख्या कीजिए तथा उसके अंगों का विस्तृत वर्णन कीजिए। \pts{10}

\medskip\rule{\linewidth}{0.2pt}\medskip

\noindent \textbf{Question 2.} Analyse the role of republican traditions in Ancient India and assess their political significance.\\[1pt]
प्राचीन भारत में गणराज्य परंपराओं की भूमिका का विश्लेषण कीजिए तथा उनके राजनीतिक महत्त्व का मूल्यांकन कीजिए। \pts{10}

\begin{center}
  \textbf{OR / अथवा}
\end{center}

\noindent Examine the system of inter-state relations and diplomacy in Ancient India with reference to Mandala theory and espionage.\\[1pt]
मंडल सिद्धान्त एवं गुप्तचर व्यवस्था के संदर्भ में प्राचीन भारत की अन्तर-राज्यीय संबंध एवं कूटनीति का परीक्षण कीजिए। \pts{10}

\medskip\rule{\linewidth}{0.2pt}\medskip

\noindent \textbf{Question 3.} Discuss the administrative system of the Mauryas with special reference to central and provincial administration.\\[1pt]
मौर्यों की प्रशासनिक व्यवस्था का वर्णन कीजिए, विशेष रूप से केंद्रीय एवं प्रांतीय प्रशासन पर बल देते हुए। \pts{10}

\begin{center}
  \textbf{OR / अथवा}
\end{center}

\noindent Discuss the administrative features of the Chola dynasty with special emphasis on local self-government.\\[1pt]
चोल वंश की प्रशासनिक विशेषताओं का वर्णन कीजिए, विशेष रूप से स्थानीय स्वशासन पर बल देते हुए। \pts{10}

\bigskip
\rule{\linewidth}{0.2pt}
\bigskip

\begin{center}
  {\large\bfseries SECTION B / खण्ड-ब}\\[2pt]
  {\normalsize (Short Answer Type Questions / लघु उत्तरीय प्रश्न)}
\end{center}

\noindent \textit{Note: Attempt all questions. Answer each question in approximately 125-150 words. Each question carries 05 marks.}\\[2pt]
\noindent \textit{नोट: सभी प्रश्नों के उत्तर दीजिए। प्रत्येक प्रश्न का उत्तर लगभग 125-150 शब्दों में दीजिए। प्रत्येक प्रश्न 05 अंकों का है।}

\bigskip

\noindent \textbf{Question 4.} Explain the nature and scope of the State in Ancient Indian thought.\\[1pt]
प्राचीन भारतीय चिंतन में राज्य के स्वरूप एवं क्षेत्र को स्पष्ट कीजिए। \pts{5}

\medskip\rule{\linewidth}{0.2pt}\medskip

\noindent \textbf{Question 5.} Discuss the origin and development of the concept of Kingship in Ancient India.\\[1pt]
प्राचीन भारत में राजत्व की अवधारणा के उद्भव एवं विकास की चर्चा कीजिए। \pts{5}

\medskip\rule{\linewidth}{0.2pt}\medskip

\noindent \textbf{Question 6.} Examine the composition and functions of the Sabha and Samiti in Ancient India.\\[1pt]
प्राचीन भारत में सभा एवं समिति के गठन एवं कार्यों का परीक्षण कीजिए। \pts{5}

\medskip\rule{\linewidth}{0.2pt}\medskip

\noindent \textbf{Question 7.} Write a short note on the administrative features of the Gupta dynasty.\\[1pt]
गुप्त वंश की प्रशासनिक विशेषताओं पर संक्षिप्त टिप्पणी लिखिए। \pts{5}

\begin{center}
  \textbf{OR / अथवा}
\end{center}

\noindent Discuss the administrative system of the Rashtrakutas in brief.\\[1pt]
राष्ट्रकूटों की प्रशासनिक व्यवस्था का संक्षेप में वर्णन कीजिए। \pts{5}

\bigskip
\rule{\linewidth}{0.2pt}
\bigskip

\begin{center}
  {\large\bfseries SECTION C / खण्ड-स}\\[2pt]
  {\normalsize (Very Short Answer Type Questions / अति लघु उत्तरीय प्रश्न)}
\end{center}

\noindent \textit{Note: Write very short notes on the following, each in about 15-20 words. Each question carries 02 marks.}\\[2pt]
\noindent \textit{नोट: निम्नलिखित प्रत्येक प्रश्न पर लगभग 15-20 शब्दों में संक्षिप्त टिप्पणी लिखिए। प्रत्येक प्रश्न 02 अंकों का है।}

\bigskip

\noindent \textbf{Question 8.}
\begin{parts}
  \item What is meant by the State?\\[1pt]
        राज्य से क्या अभिप्राय है? \pts{2}

  \item Explain the meaning of the term `Yogakshema'.\\[1pt]
        `योगक्षेम' का अर्थ स्पष्ट कीजिए। \pts{2}

  \item Identify the source of the Sanskrit passage \textit{``Svamyamatyajanapadadurgakosadandamitraprakrtayah''}.\\[1pt]
        \textit{``स्वाम्यमात्यजनपददुर्गकोशदण्डमित्राणि प्रकृतयः''} इस संस्कृत वाक्यांश का स्रोत बताइए। \pts{2}

  \item Name the important elements (limbs) mentioned in the Rigveda.\\[1pt]
        ऋग्वेद में उल्लिखित महत्वपूर्ण तत्वों/अंगों के नाम बताइए। \pts{2}

  \item State the source of the verse: \textit{``Prajasukhe sukham rajnah prajanam ca hite hitam''}.\\[1pt]
        \textit{``प्रजासुखे सुखं राज्ञः प्रजानां च हिते हितम्''} इस श्लोक का स्रोत बताइए। \pts{2}

  \item What is meant by Vairajya? Briefly explain.\\[1pt]
        `वैराज्य' से क्या अभिप्राय है? संक्षेप में स्पष्ट कीजिए। \pts{2}

  \item What are the types of victory described in Ancient Indian texts?\\[1pt]
        प्राचीन भारतीय ग्रंथों में वर्णित विजय के प्रकार कौन-कौन-से हैं? \pts{2}

  \item What is Matsyanyaya?\\[1pt]
        मात्स्यन्याय क्या है? \pts{2}

  \item What is Danda in Ancient Indian polity?\\[1pt]
        प्राचीन भारतीय राजनीति में `दण्ड' क्या है? \pts{2}

  \item Define Purushartha.\\[1pt]
        पुरुषार्थ को परिभाषित कीजिए। \pts{2}
\end{parts}

\vfill
\rule{\linewidth}{0.4pt}

\begin{center}\small
\href{https://bhu-exam.vercel.app/}{Click here to visit our website}\\[4pt]
\href{https://me-aryan.vercel.app/}{Click here to visit our contributor}
\end{center}

\end{document}
